\section{Diseño del modelo: arquitectura multietapa que unifica percepción, planificación y ejecución}

\subsection{Deficiencias clave de los GUI Agent actuales}
El expositor resumió tres tipos de problemas de los sistemas actuales: grounding visual inestable, planificación de largo alcance débil y mala generalización entre plataformas. Muchos modelos funcionan en tareas cortas, pero se degradan rápidamente en tareas de alta complejidad.

\begin{importantbox}{Tres tipos de carencias de capacidad}
\begin{enumerate}
    \item \textbf{Grounding}: percibe, pero no ubica con precisión, o hay ambigüedad al resolver el objetivo;
    \item \textbf{Planning}: ejecuta pasos individuales, pero carece de gestión global de subobjetivos;
    \item \textbf{Generalization}: el rendimiento cae abruptamente al cambiar de interfaz o de plantilla de tarea.
\end{enumerate}
\end{importantbox}

\subsection{El enfoque integrador que propone el curso}
Según los subtítulos, el expositor subraya que, mediante un paradigma generativo unificado de estilo AR, la comprensión visual, la cadena de razonamiento (monologue/reasoning trace) y la generación de acciones se incorporan a una misma interfaz de entrenamiento. Esto reduce las pérdidas de interfaz que produce el ensamblaje de múltiples modelos.

\begin{figure}[H]
\centering
\includegraphics[width=0.88\textwidth]{frame-07-model-architecture.jpg}
\caption{Video aproximadamente en 01:16:05: la arquitectura del modelo enfatiza un formato unificado de entrada y salida, así como objetivos de entrenamiento por etapas}
\end{figure}

\begin{knowledgebox}{Beneficios de ingeniería del entrenamiento por etapas}
\begin{itemize}
    \item Etapa 1: refuerza el grounding visual y la comprensión semántica de la interfaz;
    \item Etapa 2: incorpora la cadena de razonamiento y señales de planificación, para mejorar la toma de decisiones de largo alcance;
    \item Etapa 3: alinea la estrategia de ejecución con la retroalimentación de recompensa, para optimizar la tasa de finalización de tareas.
\end{itemize}
\end{knowledgebox}

\subsection{La incorporación de señales de planner}
En la parte final del curso se menciona repetidamente que el planner aporta beneficios significativos en tareas difíciles. Intuitivamente, el planner no sustituye al modelo de acción, sino que ofrece una estructura intermedia interpretable, que permite al modelo reducir la exploración a ciegas en tareas largas.

\begin{warningbox}{Planner no equivale a ``redundancia adicional''}
Si la tarea en sí tiene una cadena larga y un estado complejo, la ausencia de planner suele traducirse en más acciones inválidas y retrocesos. El costo adicional de la planificación suele ser menor que el costo de una ejecución errónea. La decisión de usar o no un planner debe tomarse por niveles, según la complejidad de la tarea.
\end{warningbox}

\subsection{Resumen de la sección}
Lo determinante en la arquitectura del modelo no es el tamaño de los parámetros, sino la combinación de ``interfaz unificada + entrenamiento por etapas + señales de planificación''. Estos tres elementos determinan en conjunto la estabilidad del sistema en tareas de alta complejidad.

